\chapter{Logistic Regression and Classification}

\section{Probabilistic Interpretation}

In contrast to regression, where outputs are real-valued, classification problems involve predicting discrete labels. A natural approach is to model the \emph{probability} of each class given the input.

\medskip

\noindent
\textbf{Binary classification.}  
Let $\mathcal{Y} = \{0,1\}$. Instead of predicting labels directly, we model
\[
P(y = 1 \mid x).
\]

\medskip

\noindent
\textbf{Logistic model.}  
We assume
\[
P(y = 1 \mid x) = \sigma(\langle w, x \rangle),
\]
where $\sigma$ is the sigmoid function
\[
\sigma(z) = \frac{1}{1 + e^{-z}}.
\]

\medskip

\noindent
\textbf{Properties of the sigmoid.}
\begin{itemize}
    \item Maps $\mathbb{R}$ to $(0,1)$
    \item Smooth and monotonic
    \item Interpretable as probability
\end{itemize}

\medskip

\noindent
\textbf{Log-odds interpretation.}  
The model implies
\[
\log \frac{P(y=1 \mid x)}{P(y=0 \mid x)} = \langle w, x \rangle.
\]

\medskip

\noindent
Thus, logistic regression models the log-odds as a linear function of the input.

\medskip

\noindent
\textbf{Multiclass extension.}  
For $K$ classes, we use the softmax function:
\[
P(y = k \mid x) = \frac{\exp(\langle w_k, x \rangle)}{\sum_{j=1}^K \exp(\langle w_j, x \rangle)}.
\]

\medskip

\noindent
This formulation ensures that probabilities are positive and sum to one.

\section{Cross-Entropy Loss}

To train the model, we need a loss function that reflects probabilistic accuracy.

\medskip

\noindent
\textbf{Binary cross-entropy loss.}
\[
\ell(f(x), y) = -\big( y \log p(x) + (1-y)\log(1 - p(x)) \big),
\]
where $p(x) = P(y=1 \mid x)$.

\medskip

\noindent
\textbf{Multiclass cross-entropy.}
\[
\ell(f(x), y) = -\log P(y \mid x).
\]

\medskip

\noindent
\textbf{Interpretation.}
\begin{itemize}
    \item Penalizes confident but incorrect predictions heavily
    \item Encourages calibrated probabilities
\end{itemize}

\medskip

\noindent
\textbf{Geometric insight.}  
Cross-entropy measures the discrepancy between the predicted distribution and the true label distribution.

\medskip

\noindent
\textbf{Practical advantage.}  
Unlike squared loss, cross-entropy leads to well-behaved gradients for classification tasks.

\section{Maximum Likelihood Perspective}

Logistic regression can be derived from the principle of maximum likelihood.

\medskip

\noindent
\textbf{Likelihood function.}
\[
L(w) = \prod_{i=1}^n P(y_i \mid x_i; w).
\]

\medskip

\noindent
\textbf{Log-likelihood.}
\[
\ell(w) = \sum_{i=1}^n \log P(y_i \mid x_i; w).
\]

\medskip

\noindent
Maximizing the log-likelihood is equivalent to minimizing the empirical cross-entropy loss:
\[
\min_w \; -\sum_{i=1}^n \log P(y_i \mid x_i; w).
\]

\medskip

\noindent
\textbf{Binary case.}  
Substituting the sigmoid model yields a convex optimization problem.

\medskip

\noindent
\textbf{Gradient.}  
For binary classification,
\[
\nabla_w \ell(w) = \sum_{i=1}^n ( \sigma(\langle w, x_i \rangle) - y_i ) x_i.
\]

\medskip

\noindent
\textbf{Insight.}
\begin{itemize}
    \item The gradient reflects prediction error weighted by input features
    \item Training adjusts parameters to reduce discrepancies between predicted probabilities and labels
\end{itemize}

\medskip

\noindent
\textbf{Regularization.}  
As in linear regression, one often adds an $\ell_2$ penalty:
\[
\min_w \; -\sum_{i=1}^n \log P(y_i \mid x_i; w) + \lambda \|w\|_2^2.
\]

\medskip

\noindent
\textbf{Interpretation.}  
This corresponds to maximum a posteriori (MAP) estimation under a Gaussian prior.

\section{Decision Boundaries}

Although logistic regression models probabilities, it ultimately produces decisions.

\medskip

\noindent
\textbf{Binary decision rule.}
\[
\hat{y} =
\begin{cases}
1 & \text{if } P(y=1 \mid x) \geq \tfrac{1}{2}, \\
0 & \text{otherwise}.
\end{cases}
\]

\medskip

\noindent
\textbf{Equivalent condition.}
\[
\langle w, x \rangle \geq 0.
\]

\medskip

\noindent
\textbf{Decision boundary.}  
The set of points satisfying
\[
\langle w, x \rangle = 0
\]
defines a hyperplane in $\mathbb{R}^d$.

\medskip

\noindent
\textbf{Geometric interpretation.}
\begin{itemize}
    \item The input space is divided into two half-spaces
    \item The normal vector $w$ determines orientation
\end{itemize}

\medskip

\noindent
\textbf{Multiclass case.}  
Decision boundaries are determined by comparing linear scores:
\[
\langle w_k, x \rangle = \langle w_j, x \rangle.
\]

\medskip

\noindent
\textbf{Limitation.}  
Logistic regression produces \emph{linear decision boundaries}. It cannot capture nonlinear structure without feature transformations.

\medskip

\noindent
\textbf{Extension.}  
Nonlinear decision boundaries can be obtained by:
\begin{itemize}
    \item feature mappings,
    \item kernel methods,
    \item neural networks.
\end{itemize}

\section{A Unifying Perspective}

Logistic regression connects several key ideas:

\begin{itemize}
    \item \textbf{Probabilistic modeling:} predicting conditional probabilities
    \item \textbf{Optimization:} minimizing cross-entropy loss
    \item \textbf{Geometry:} linear decision boundaries
\end{itemize}

\medskip

\noindent
\textbf{Comparison with linear regression.}
\begin{itemize}
    \item Linear regression predicts values
    \item Logistic regression predicts probabilities
\end{itemize}

\medskip

\noindent
\textbf{Key insight.}  
Classification can be viewed as modeling uncertainty about discrete outcomes.

\section{Outlook}

This chapter introduced logistic regression as a fundamental classification model:
\begin{itemize}
    \item probabilistic interpretation,
    \item cross-entropy loss,
    \item maximum likelihood estimation,
    \item geometric structure of decision boundaries.
\end{itemize}

\medskip

In the next chapter, we extend linear models using feature maps and kernels, enabling them to capture nonlinear relationships.

\medskip

\noindent
\textbf{Guiding principle.}  
Effective classification requires both expressive models and appropriate probabilistic formulations.